\section{例题与完整推导}

\subsection{统一计算框架}

对三维网格尺寸
\[
(G_x,G_y,G_z)
\]
与三维块尺寸
\[
(B_x,B_y,B_z),
\]
依次计算
\[
\begin{aligned}
N_{\mathrm{threads/block}} &= B_xB_yB_z,\\
N_{\mathrm{blocks}} &= G_xG_yG_z,\\
N_{\mathrm{threads,total}} &=
N_{\mathrm{threads/block}}N_{\mathrm{blocks}}.
\end{aligned}
\]
未使用的维度取 1，不能把它们从概念上混成 0；尺寸为 0 会意味着不启动相应工作，而一维或二维配置的剩余维度实际为 1。

\subsection{例题一：一维网格与二维线程块}

\begin{examplebox}[title=题目]
一个 CUDA 核函数启动 32 个线程块，每个线程块包含 $32\times32$ 个线程。求每块线程数与总启动线程数。
\end{examplebox}

网格可写为
\[
\texttt{gridDim}=(32,1,1),
\]
线程块可写为
\[
\texttt{blockDim}=(32,32,1).
\]
因此每块线程数为
\[
N_{\mathrm{threads/block}}
=32\times32\times1
=1024.
\]
网格中的块数为
\[
N_{\mathrm{blocks}}
=32\times1\times1
=32.
\]
总线程数为
\[
N_{\mathrm{threads,total}}
=1024\times32
=32768.
\]

\begin{keybox}[title=答案]
每个线程块包含 $1024$ 个线程；整个网格共启动 $32768$ 个线程。
\end{keybox}

\subsection{例题二：二维网格与二维线程块}

\begin{examplebox}[title=题目]
线程块配置为 $16\times16$，网格尺寸为 $32\times16$ 个块。求每块线程数、块总数与整个网格的线程总数。
\end{examplebox}

对应尺寸为
\[
\texttt{blockDim}=(16,16,1),
\qquad
\texttt{gridDim}=(32,16,1).
\]
每块线程数为
\[
N_{\mathrm{threads/block}}
=16\times16\times1
=256.
\]
块总数为
\[
N_{\mathrm{blocks}}
=32\times16\times1
=512.
\]
因此总线程数为
\[
N_{\mathrm{threads,total}}
=256\times512
=131072.
\]
也可一次性验证：
\[
(32\times16)\times(16\times16)=131072.
\]

\begin{keybox}[title=答案]
每块包含 $256$ 个线程；网格含 $512$ 个块；整个网格共启动 $131072$ 个线程。
\end{keybox}

\subsection{例题三：三维网格与三维线程块}

\begin{examplebox}[title=题目]
网格尺寸为 $8\times4\times2$ 个块，每个线程块尺寸为 $4\times8\times8$。求每块线程数与整个网格的线程总数。
\end{examplebox}

每块线程数为
\[
N_{\mathrm{threads/block}}
=4\times8\times8
=256.
\]
网格中的块数为
\[
N_{\mathrm{blocks}}
=8\times4\times2
=64.
\]
总线程数为
\[
N_{\mathrm{threads,total}}
=256\times64
=16384.
\]

\begin{keybox}[title=答案]
每块包含 $256$ 个线程；网格含 $64$ 个块；整个网格共启动 $16384$ 个线程。
\end{keybox}

\subsection{综合例题：最后一个有效像素与第一个无效线程}

考虑 $76\times62$ 图像，块尺寸为 $16\times16$，因此
\[
\texttt{gridDim}=(5,4,1).
\]
图像最后一个有效像素坐标为
\[
(r,c)=(61,75).
\]
取最后一个线程块
\[
\texttt{blockIdx}=(4,3),
\]
并取块内线程
\[
\texttt{threadIdx}=(11,13).
\]
其坐标为
\[
\begin{aligned}
c&=4\times16+11=75,\\
r&=3\times16+13=61.
\end{aligned}
\]
因此该线程恰好负责最后一个有效像素，线性下标为
\[
p=61\times76+75=4711,
\]
与长度 $76\times62=4712$ 的数组最后下标一致。

若只把块内 $x$ 坐标增加 1，即取 \texttt{threadIdx.x}=12，则
\[
c=4\times16+12=76.
\]
此时 $c<76$ 不成立，该线程是最靠近右边界的无效线程之一。这个例子说明，线程的块内索引本身完全合法，并不意味着它对应的数据坐标合法；数据边界必须单独判断。

\subsection{常见计算错误}

\begin{itemize}
  \item 把“网格中的块数”与“每块线程数”混为同一个量；
  \item 只乘二维图中显式写出的分量，却忘记未使用维度应乘以 1；
  \item 把 \texttt{gridDim} 当成线程总数，或把 \texttt{blockDim} 当成块总数；
  \item 认为启动线程总数必然等于有效数据元素数，忽略向上取整产生的边界线程；
  \item 在二维索引中把 $x$ 当作行、$y$ 当作列，导致宽度与高度的判断互换。
\end{itemize}
